From the analyse of the results of the measurements in Appedix XX, it is shown, that the height of the antennas and the distance between them have a much higher influence on the PL, than the other parameter, which is the placement, antenna type and polarization. As these parameters have much less influence on the PL, there will only be looking at the different height combinations and distance between the antennas.

Firstly there will be looked at the case, where the height of the transmitting antenna and the receiving antenna is at the same heights, over the different distance.

\begin{figure}
\centering
\includegraphics[width=0.3\textwidth]{figures/bleach-3712665.jpg}
\caption{Mesurements of PL for transmitter and receiver antenna at the same height.}
\label{Meas1}
\end{figure}

It is seen on figure \ref{Meas1}, that the longer the distance, the higher the PL. This comes from the same principle, that the Friss model is based on. When looking at the graph compared to changing the heights, it is seen, that the lower the antennas is placed, the higher PL. This comes from the principle, that reflections comes from the ground and comes with (in most cases) destructive interference, which the Two Ray model is based on. 

Secondly there will be looked at the case, where the height of the transmitting antenna is kept at a single height, both low at 0.01m and high at 2.00m, while the height of the receiving antenna is varying.

\begin{figure}
\centering
\includegraphics[width=0.3\textwidth]{figures/bleach-3712665.jpg}
\caption{Mesurements of PL for transmitter at 0.01m and receiver antenna at varying heights}
\label{Meas1}
\end{figure}

\begin{figure}
\centering
\includegraphics[width=0.3\textwidth]{figures/bleach-3712665.jpg}
\caption{Mesurements of PL for transmitter at 2.00m and receiver antenna at varying heights}
\label{Meas1}
\end{figure}